% TOP_PROBLEM: {"id": 156, "title": "Restricted Sumsets in Partitions", "category_id": 2, "category": {"id": 2, "name": "combinatorics", "display_name": "Combinatorics", "description": "Counting problems, graph theory, discrete structures.", "slug": "combinatorics", "order_index": 2, "created_at": "2026-07-31T15:26:25.671Z"}, "status": "open"}
% FIRST_POSTED: null
% ENTRY_KIND: statement_only
% Imported from: batch03.tex; UnsolvedMath ID 156
\documentclass[11pt]{article}
\usepackage[margin=25mm]{geometry}
\usepackage{fontspec}
\setmainfont{Latin Modern Roman}
\usepackage{amsmath,amssymb,mathtools}
\usepackage{unicode-math}
\setmathfont{Latin Modern Math}
\usepackage{xurl,hyperref}
\hypersetup{colorlinks=true,urlcolor=blue}
\setcounter{secnumdepth}{0}
\setlength{\parindent}{0pt}
\setlength{\parskip}{5pt}
\setlength{\emergencystretch}{3em}
\newcommand{\sref}[2]{\hyperlink{source:#1:#2}{[S#2]}}
\newcommand{\cataloguescope}[1]{\paragraph{Recorded target.}#1}
\begin{document}
% UnsolvedMath ID 156; source catalogue code GREEN-068
\setcounter{section}{155}
\section{Restricted Sumsets in Partitions}\label{156}
{\small\sffamily UnsolvedMath ID 156 · Combinatorics}
\subsection{Definitions and mathematical statement}

\paragraph{Shared notation.}$\mathbb N=\{1,2,\ldots\}$, and $\mathbb Z,\mathbb Q,\mathbb R,\mathbb C$ denote the integers, rationals, reals and complex numbers. Any other conventions are specified locally.


For a set $B\subset\mathbb Z$, its restricted sumset is $B\widehat{+}B=\{b+b':b,b'\in B,\ b\ne b'\}$. For integers $N\geq2$ and $1\leq k\leq N$, define
\[
 R_k(N)=\min_{c:[N]\to[k]}
 \bigl|\{a+b:1\leq a<b\leq N,\ c(a)=c(b)\}\bigr|.
\]
Empty color classes are allowed. The counted set is the union of the restricted sumsets of the color classes; a sum attained in several classes is counted once.
\paragraph{Statement recorded in the supplied source.}
Determine, as a function of $N$, the values of $k$ for which every $k$-coloring of $[N]$ satisfies $R_k(N)\geq N/10$. Equivalently, determine the growth of the threshold
\[
 \kappa(N)=\max\{k\in\{1,\ldots,N\}:R_k(N)\geq N/10\}.
\]
The number of colors may grow with $N$; this is not restricted to the already-known case of fixed $k$. \sref{156}{2}
\subsection{Short English statement}
How many colors can an interval be partitioned into while guaranteeing that the union of all same-color sums of two distinct numbers has size at least one tenth of the interval?
\subsection{Sources}
\begin{description}
\item[\hypertarget{source:156:1}{S1}] ULAM.AI and the UnsolvedMath Contributors, \emph{UnsolvedMath: A Curated Collection of Open Mathematics Problems}, record 156 (GREEN-068). CC BY 4.0. \url{https://huggingface.co/datasets/ulamai/UnsolvedMath}.
\item[\hypertarget{source:156:2}{S2}] Ben Green, \emph{100 Open Problems}, maintained author collection, Problem 25 and its comments, accessed 7 October 2026. \url{https://people.maths.ox.ac.uk/greenbj/papers/open-problems.pdf}.
\end{description}
\label{end:156}
\end{document}
